\begin{table}[H]
\centering
\caption{\textbf{For-profit--public differences in early-term birth and newborn outcomes}}
\label{tab:health}
\small\setlength{\tabcolsep}{5pt}
\resizebox{\ifdim\width>\linewidth \linewidth\else\width\fi}{!}{%
\begin{tabular}{lcccc}
\toprule
 & (1) & (2) & (3) & (4) \\
 & \multicolumn{2}{c}{Early-term birth} & Low birthweight & Five-minute \\
 & \multicolumn{2}{c}{(37--38 weeks)} & ($<$2500g) & Apgar $<$ 7 \\
\cmidrule(lr){2-3}\cmidrule(lr){4-4}\cmidrule(lr){5-5}
For-profit establishment & 14.14$^{***}$ & 11.73$^{***}$ & -2.99$^{***}$ & -0.53$^{***}$ \\
 & (0.74) & (0.72) & (0.66) & (0.07) \\
\addlinespace[2pt]
\midrule
Maternal controls & No & Yes & Yes & Yes \\
Municipality fixed effects & Yes & Yes & Yes & Yes \\
Year fixed effects & Yes & Yes & Yes & Yes \\
Observations & 23,141,498 & 22,019,981 & 22,436,582 & 22,144,638 \\
\bottomrule
\end{tabular}}
\begin{minipage}{\linewidth}\footnotesize
\textit{Notes:} Birth-level regressions, SINASC 2010--2024, for-profit vs
public establishments; coefficients in percentage points. Maternal controls:
age, age$^2$, education, race. Mothers differ across sectors, so the
early-term coefficient is an associational difference between establishment
sectors and is not interpreted as the causal effect of prelabor scheduling.
Standard errors, clustered by municipality, are reported in parentheses.
\newline
\footnotesize Significance levels: *** p$<$0.01, ** p$<$0.05, * p$<$0.10.
\end{minipage}
\end{table}
